\documentclass[instructions.tex]{subfiles}
\begin{document}
 
\chapter{Perdre le temps, c'est être insensé}
 

\primo Lorsque j'étais enfant, dit saint Paul, je pensais comme un enfant, je parlais comme un enfant. Combien de gens qui se piquent d'avoir de l'esprit, vivent en enfant, agissent toute leur vie avec moins de jugement et ne réfléchissent pas plus que des enfans! 

Que penserait-on d'un homme pauvre qui, ayant une somme d'argent, s'occuperait à la perdre, à la faire voltiger sur l'eau, comme un enfant, s'il alléguait pour raison qu'il le fait pour passer le temps ? que devez-vous penser de vous-même qui, étant si pauvre en mérites, perdez un temps avec lequel vous pourriez vous enrichir pour l'éternité ? 
 
 Le monde regarderait comme une folie de mépriser l'occasion de faire sa fortune; mais c'est une bien plus grande folie de perdre à tout moment l'occasion de gagner le ciel. Si un roi vous ordonnait de prendre, pendant une heure, dans ses trésors telle somme que vous voudriez pour payer vos dettes et pour acheter un riche domaine, avec cette condition que, si vous ne prenez pas toute la somme nécessaire, vous serez condamné à mort, pour avoir perdu cette occasion de satisfaire vos créanciers et d'établir votre famille; ne seriez-vous pas un insensé, si, au lieu de profiter de cette conjoncture, vous vous amusiez à vous divertir ou à cueillir des fleurs ?

 Votre conduite est encore plus insensée lorsque vous perdez le temps. Qu'est ce que le temps de la vie en comparaison de l'éternité ? C'est moins qu'un jour, c'est moins qu'une heure. Pendant ce temps si court, Dieu vous commande de travailler pour satisfaire à sa justice et pour vous enrichir de mérites; de telle sorte que, si à votre mort il vous manque quelque chose de toute ce qui est nécessaire pour entrer dans le ciel, vous serez perdu sans ressource. 

 Comprenez donc à quoi vous vous exposez lorsque vous employez ce temps à la vanité, à des plaisirs indignes de vous, à acquérir des biens dont vous ne jouirez que deux jours. Si vous prétendez excuser une telle conduite, vous n'êtes pas aveugle, mais extravagant. 

 Que penserait-on d'un homme qui emploierait tous ses biens à faire creuser jusqu'aux entrailles de la terre pour en tirer une paille d'or ? Que devez-vous pensez de vous-même qui vous consumez en soins pour les choses de la terre, qui n'employez à tant d'inutilités un temps dont les momens sont si précieux, que tout l'or du monde ne peut vous rendre ce qu'un moment perdu vous ôte, et qu'un moment de temps ne peut s'acheter avec toutes les richesses de l'univers. 

\secundo Un jour viendra que vous verserez des larmes, en souhaitant quelques momens de tant d'heures que vous perdez, sans pouvoir en obtenir un seul. Saint Grégoire rapporte qu'un homme qui avait passé sa vie dans le crime et dans l'oubli de Dieu, se trouvant, la nuit, surpris d'une maladie mortelle, s'écria en expirant:\textit{Ah! grand Dieu, accordez-moi quelques momens de vie. Seigneur, la vie jusqu'à demain!} Mais il ne put obtenir ce peu de temps pour mettre ordre à sa conscience, après avoir perdu tant de momens et tant d'occasions de se sauver. Profitons donc de tous les momens que Dieu nous donne; ils sont comptés, personne ne peut y ajouter un seul instant. 

Un officier de Charles-Quint,, sur le point de mourir, le pria de lui prolonger la vie de quelques jours; mais ce prince n'avait pas plus de pouvoir sur le temps que les autres hommes, qui n'ont pas un moment en leur disposition pour eux-mêmes, bien moins encore pour les autres.

Hélas! les jours du salut s'en vont, dit saint Bernard, et personne n'y pense\footnote{Transeunt dies salutis, et nemo recogitat. S.Bern.}. Si vous avez perdu le passé, pensez sérieusement à l'avenir, profitez du présent. S'il n'est rien de plus affligeant, à la mort, que d'avoir abusé du temps, il n'est rien de plus consolant que de l'avoir saintement employé.


\section{Résolutions}

\primo Je regarderai comme perdu, et je me reprocherai sincérement tout le temps employé à des choses qui ne seraient ni de mon devoir, ni des bienséances, ni d'une utilité réelle, et que je ne pourrais offrir à Dieu.

\secundo Je renonce à toute lecture dangereuse, à tout divertissement où préside le démon, à toute visite qui ne pourrait avoir pour but que de me faire perdre le temps et exposer mon salut. 

\prayer{O Mon Dieu! je vous demande pardon de l'abus que j'ai fait jusqu'ici du temps. Inspirez-moi une grande idée de celui que vous voulez bien m'accorder encore, et faites que je l'emploie si bien dans la suite, que le ciel soit un jour la récompense de l'usage que j'en aurai désormais fait.}

\end{document}
